We present \LRL\ (\lrl), a language and compiler that combines three mechanisms from different traditions: dependent types checked by a separate trusted kernel in the style of Lean and Coq; Rust's ownership discipline, with values that may be used at most once, closures classified as \lstinline|Fn|, \lstinline|FnMut| or \lstinline|FnOnce|, and borrows checked with non-lexical lifetimes; and hygienic Lisp-style macros.
We make three contributions, each backed by the implementation.
First, we define function kinds precisely---a kind says how calling a function uses the environment it captured---and give the typing and ownership rules of the kernel. Kinds interact with dependent elimination, the type-theoretic form of pattern matching and structural recursion: the function that implements it, the recursor, calls each case a number of times that depends on the data, so a case that may run repeatedly must not consume what it captures. The kernel enforces this, and for a non-dependent core calculus with function kinds and an iterator we mechanize in Lean a proof that well-owned programs never consume a resource twice.
Second, we state which component guarantees what: the kernel, which is the logical trusted base, guarantees that no owned value is moved twice or used directly after a move; a separate checker on the compiler's intermediate representation, outside the kernel, guarantees that borrows neither conflict nor outlive their owners; every definition, including code produced by macros, which cannot observe types or ownership, must pass both.
Third, we validate the implementation with complete case studies that combine length-indexed vectors, kernel-checked proofs and an affine protocol channel whose operations are generated by a macro; a corpus of ownership violations, each written by hand and generated by a macro; a stage-by-stage comparison of the two checkers; and the same properties tested in Lean~4, Idris~2, Rust and Racket with programs that we compiled and ran.
\lrl\ is an alpha-stage research artifact; we report its limitations together with its guarantees.
